\section{Further research questions}

\begin{enumerate}
\item \textbf{Formalize the restricted equality kernel.}
Add a new derivability predicate to the pinned Lean development and prove
the fresh-alias theorem, unrestricted crossed transport, and exact
equivalence with \texttt{NPrf}. The useful boundary is the alias lemma:
its proof needs computation and root joining, but no variable-substitution
instance. This isolates the noncircular bootstrap for a formal audit.

\item \textbf{Find the least residual equality interface.}
The useful instances of Param are root joining and substitution at a
fresh variable. Its fallback is only symmetry. Can a smaller total
operation express exactly those indispensable instances without a larger
syntax test? A proposal must account for arbitrary ground equality
assumptions, rather than showing only that computation equations survive.

\item \textbf{Control certificate expansion.}
The primitive checker becomes simpler while the proof producer builds an
alias scope around each arbitrary template transport. Quantify the
shortest expansion and compare tree certificates with explicit sharing.
A smaller kernel can be accompanied by larger proofs; both costs deserve
measurement.

\item \textbf{Obtain a substitution-stable single introduction.}
The literal-witness convention remains sensitive to numerical
substitution. The current diagonal-guard work supplies a stable alternative,
but it costs a more elaborate shape test. A smaller stable guard must also
permit recovery of case introduction; a conditional closure with an
unproved base cannot simply be used as a progressive closure.

\item \textbf{Separate arithmetic reductions from coding reductions.}
The predecessor/equality-reflection bootstrap remains the obstacle for
weaker successor-rigidity or zero-equation induction proposals. Solving it
would change the arithmetic foundation, while the present theorem changes
the equality interface. The latter does not settle the former.

\item \textbf{State a reproducible minimization problem.}
A meaningful lower bound needs a fixed public contract, source language,
normalization policy, and cost metric. The present result improves one
matched byte coordinate and restricts the equality schema; it supplies no
lower bound proving that these six explicitly split schemas are optimal.
\end{enumerate}

\section{Provenance and exact status}

The construction is a new restriction and integration of ingredients in
the supplied project. The definitions of scalar iteration, paired
closure, clause contexts, and the single witness-directed introduction
are inherited. The idea of introducing a fresh numerical alias also
appears in the optional proof in \texttt{template\_free}. The present
argument isolates the smaller equality interface it actually requires.

The archive includes unchanged inherited bridge fragments for reference:
the native choice basis, ordinary logical bridge, scalar predecessor,
equality reflection, and compilation of primitive recursion. Their role
is to document the existing arithmetic dependency. They are not counted
as newly proved results. The main paper gives the complete new reduction
and names the precise inherited Lean statements used at the last step:
\texttt{native\_iff\_pra}, \texttt{pra\_iff\_native},
\texttt{native\_roundTrip}, \texttt{pra\_native\_roundTrip}, and
\texttt{native\_inconsistent\_iff}.

Original ProveIt material is distributed under its included MIT No
Attribution license. The new checker, proof producers, proof notes, and
reproduction scripts may be reused under the same terms. The repository
was read at the pinned revision; this package does not publish changes
to it.

